% =====================================================================
\section{A neighbouring hypothesis: BDNF and serotonin}
It is not enough for a neuron to receive nourishment: to maintain its processes and synapses it needs support signals from its neighbours --- \emph{neurotrophic factors}. The principal one in the cortex is the protein BDNF. It is released by working excitatory cells and received by the TrkB receptor on neighbouring cells, PV cells among them. The \emph{neurotrophic hypothesis} of schizophrenia asserts that this support is insufficient: in development this yields faulty assembly of networks, and in the adult a reduced capacity of networks to reorganise.

\emph{Serotonin} is another transmitter which, alongside dopamine, governs the overall operating regime of the cortex. It has long been linked to psychosis: hallucinogens act on its receptors. One of these receptors governs BDNF production, so the two lines converge.
\facts
\begin{itemize}
\item In the frontal cortex of deceased patients the amount of BDNF mRNA is 28\,\% lower and that of its receptor TrkB 27\,\% lower. The more TrkB is reduced in a patient, the more GAD67 is reduced (correlation 0.74).\tP
\item In mice with the TrkB receptor reduced by 75\,\%, GAD67 and parvalbumin mRNA are 23 and 40\,\% lower. In mice with BDNF itself reduced by 80\,\% they are unchanged. The authors conclude that what matters is signalling through the receptor.\tP
\item In the blood of patients BDNF is moderately reduced (a synthesis of 41 studies, more than 7000 people). The reduction grows with the duration of the illness and is unrelated to symptom severity; whether it is cause or consequence the authors leave open.\tP
\item BDNF production depends on the cell's activity and on calcium, which enters in part through the NMDA receptor.\tU
\item The serotonin receptor 5-HT2A governs BDNF production in the frontal cortex and the hippocampus; the relation is bidirectional.\tP
\item Hallucinogens (LSD, psilocybin) act through the 5-HT2A receptor; most second-generation antipsychotics block not only the D2 dopamine receptors but this one as well.\tU
\item Antipsychotics and antidepressants raise BDNF production (per a 2005 review).\tP
\end{itemize}
\begin{note}
The neurotrophic hypothesis is not new: reviews under that name appeared in 2004--2005. It does not exclude the hypothesis of this brochure: both describe a weakening of PV cells, but place a different link first. Which of them is primary, and whether both are consequences of a common cause, has not been established. A search conducted while preparing the brochure found no single published model uniting serotonin, BDNF and oxidative stress.
\end{note}
\begin{ineq}
In terms of equations BDNF is another candidate for the role of the slow link ``network activity $\to$ support of the PV cell'': it is released by active cells and acts over hours and days. The construction is the same as in section~7: a slow variable which activity maintains and which itself maintains inhibition. The difference is that the source of support is the neighbouring excitatory cells, so the loop closes through the network rather than inside a single cell.\tD
\end{ineq}
